\section{AMC 12 12B 2005}
        \begin{enumerate}
        	\item (\textit{AMC 12 12B 2005}) A scout troop buys $1000$ candy bars at a price of five for $2$ dollars. They sell all the candy bars at the price of two for $1$ dollar.  What was their profit, in dollars?


 $\mathrm{(A)}\ 100      \qquad \mathrm{(B)}\ 200      \qquad \mathrm{(C)}\ 300      \qquad \mathrm{(D)}\ 400      \qquad \mathrm{(E)}\ 500$

 \textbf{Answer: }A



\textbf{Solution}


 \begin{align*} \mbox{Expenses} &= 1000 \cdot \frac25 = 400 \\ \mbox{Revenue}  &= 1000 \cdot \frac12 = 500 \\ \mbox{Profit}   &= \mbox{Revenue} - \mbox{Expenses} = 500-400 = \boxed{\textbf{(A) }100}. \end{align*}
Note: Revenue is a gain.

 



\textbf{Solution 2 (Faster)}

Note that the troop buys $10$ candy bars at a price of $4$ dollars and sells $10$ bars at a price of $5$ dollars. So the troop gains $1$ dollar for every $10$ bars. So therefore we divide $1000 \div 10 = 100$. So our answer is $\boxed{\textbf{(A) }100}$.~HyperVoid

 


	\item (\textit{AMC 12 12B 2005}) A positive number $x$ has the property that $x\%$ of $x$ is $4$.  What is $x$?


 $\mathrm{(A)}\ 2      \qquad \mathrm{(B)}\ 4      \qquad \mathrm{(C)}\ 10      \qquad \mathrm{(D)}\ 20      \qquad \mathrm{(E)}\ 40$

 \textbf{Answer: }D



\textbf{Solution 1}

Since $x\%$ means $0.01x$, the statement "$x\% \text{ of } x \text{ is 4}$" can be rewritten as "$0.01x \cdot x = 4$":

 $0.01x \cdot x=4 \Rightarrow x^2 = 400 \Rightarrow x = \boxed{\textbf{(D) }20}.$




\textbf{Solution 2}

Try the answer choices one by one. Upon examination, it is quite obvious that the answer is $\boxed{\textbf{(D) }20}.$ Very fast.

 
Solution by franzliszt

 


	\item (\textit{AMC 12 12B 2005}) Brianna is using part of the money she earned on her weekend job to buy several equally-priced CDs. She used one fifth of her money to buy one third of the CDs. What fraction of her money will she have left after she buys all the CDs? 


 $\mathrm{(A)}\ \frac15      \qquad \mathrm{(B)}\ \frac13      \qquad \mathrm{(C)}\ \frac25      \qquad \mathrm{(D)}\ \frac23      \qquad \mathrm{(E)}\ \frac45$

 \textbf{Answer: }C



\textbf{Solution}

Let $m =$ Brianna's money.  We have $\frac15 m = \frac13 (\mbox{CDs}) \Rightarrow \frac35 m = (\mbox{CDs})$.  Thus, the money left over is $m-\frac35m = \frac25m$, so the answer is $\boxed{\textbf{(C) }\frac{2}{5}}$.

 This was just a simple manipulation of the equation. No solving was needed!

 


	\item (\textit{AMC 12 12B 2005}) At the beginning of the school year, Lisa's goal was to earn an A on at least $80\%$ of her $50$ quizzes for the year.  She earned an A on $22$ of the first $30$ quizzes.  If she is to achieve her goal, on at most how many of the remaining quizzes can she earn a grade lower than an A?


 $\mathrm{(A)}\ 1      \qquad \mathrm{(B)}\ 2      \qquad \mathrm{(C)}\ 3      \qquad \mathrm{(D)}\ 4      \qquad \mathrm{(E)}\ 5$

 \textbf{Answer: }B



\textbf{Solution}

Lisa's goal was to get an $A$ on $80\% \cdot 50 = 40$ quizzes.  She already has $A$'s on $22$ quizzes, so she needs to get $A$'s on $40-22=18$ more.  There are $50-30=20$ quizzes left, so she can afford to get less than an $A$ on $20-18=\boxed{\textbf{(B) }2}$ of them.

 Here, only the $A$'s matter... No complicated stuff!

 


	\item (\textit{AMC 12 12B 2005}) An $8$-foot by $10$-foot floor is tiled with square tiles of size $1$ foot by $1$ foot.  Each tile has a pattern consisting of four white quarter circles of radius $1/2$ foot centered at each corner of the tile.  The remaining portion of the tile is shaded.  How many square feet of the floor are shaded?


 
                    \begin{figure}[H]
                        \centering
                        \includegraphics[width=0.5\textwidth]{Problem/AMC_12_12B_2005_5_Problem_Figure_1.png}
                    \end{figure}
                    
$\mathrm{(A)}\ 80-20\pi      \qquad \mathrm{(B)}\ 60-10\pi      \qquad \mathrm{(C)}\ 80-10\pi      \qquad \mathrm{(D)}\ 60+10\pi      \qquad \mathrm{(E)}\ 80+10\pi$

 \textbf{Answer: }A



\textbf{Solution}

There are $80$ tiles.  Each tile has $[\mbox{square} - 4 \cdot (\mbox{quarter circle})]$ shaded.  Thus:

 
 \begin{align*} \mbox{shaded area} &= 80 \left( 1 - 4 \cdot \dfrac{1}{4} \cdot \pi \cdot \left(\dfrac{1}{2}\right)^2\right) \\ &= 80\left(1-\dfrac{1}{4}\pi\right) \\ &= \boxed{\textbf{(A) }80-20\pi}. \end{align*}
4 quarters of a circle is a circle so that may save you 0.5 seconds :)

 


	\item (\textit{AMC 12 12B 2005}) In $\triangle ABC$, we have $AC=BC=7$ and $AB=2$. Suppose that $D$ is a point on line $AB$ such that $B$ lies between $A$ and $D$ and $CD=8$. What is $BD$?


 $\mathrm{(A)}\ 3      \qquad \mathrm{(B)}\ 2\sqrt{3}      \qquad \mathrm{(C)}\ 4      \qquad \mathrm{(D)}\ 5      \qquad \mathrm{(E)}\ 4\sqrt{2}$

 \textbf{Answer: }A



\textbf{Solution 1}

Draw height $CH$ (Perpendicular line from point C to line AD). We have that $BH=1$. By the \href{https://artofproblemsolving.com/wiki/index.php?title=Pythagorean\_Theorem}{Pythagorean Theorem}, $CH=\sqrt{48}$. Since $CD=8$, $HD=\sqrt{8^2-48}=\sqrt{16}=4$, and $BD=HD-1$, so $BD=\boxed{\textbf{(A) }3}$.

 



\textbf{Solution 2 (Trig)}

After drawing out a diagram, let $\angle{ABC}=\theta$. By the Law of Cosines, $7^2=2^2+7^2-2(7)(2)\cos{\theta} \rightarrow 0=4-28\cos{\theta} \rightarrow \cos{\theta}=\frac{1}{7}$. In $\triangle CBD$, we have $\angle{CBD}=(180-\theta)$, and using the identity $\cos(180-\theta)=-\cos{\theta}$ and Law of Cosines one more time: $8^2=7^2+x^2-2(7)(x)\left( \frac{-1}{7} \right) \rightarrow 64=49+x^2+2x \rightarrow x^2+2x-15=0$. The only positive value for $x$ is $3$, which gives the length of $\overline{BD}$. Thus the answer is $\boxed{\textbf{(A) }3}$.

 



\textbf{Solution 3 (Stewart's Theorem)}

Let $BD=k$. Then, by \href{https://artofproblemsolving.com/wiki/index.php?title=Stewart\%27s\_Theorem}{Stewart's Theorem}, 

 $2k(2+k)+7^2(2+k)=7^2k+8^2\cdot 2 \implies k^2+2k-15=0 \implies k=\boxed{\textbf{(A) }3}$



	\item (\textit{AMC 12 12B 2005}) What is the area enclosed by the graph of $|3x|+|4y|=12$?


 $\mathrm{(A)}\ 6      \qquad \mathrm{(B)}\ 12      \qquad \mathrm{(C)}\ 16      \qquad \mathrm{(D)}\ 24      \qquad \mathrm{(E)}\ 25$

 \textbf{Answer: }D



\textbf{Solution 1}

If we get rid of the absolute values, we are left with the following 4 equations (using the logic that if $|a|=b$, then $a$ is either $b$ or $-b$):

 
 \begin{align*} 3x+4y=12 \\ -3x+4y=12 \\ 3x-4y=12 \\ -3x-4y=12 \end{align*}


 We can then put these equations in slope-intercept form in order to graph them.

 
 \begin{align*} 3x+4y=12 \,\implies\, y=-\dfrac{3}{4}x+3\\ -3x+4y=12\,\implies\, y=\dfrac{3}{4}x+3\\ 3x-4y=12\,\implies\, y=\dfrac{3}{4}x-3\\ -3x-4y=12\,\implies\, y=-\dfrac{3}{4}x-3\end{align*}


 Now you can graph the lines to find the shape of the graph:

 
                    \begin{figure}[H]
                        \centering
                        \includegraphics[width=0.5\textwidth]{Solution/AMC_12_12B_2005_7_Solution_Figure_1.png}
                    \end{figure}
                    

 We can easily see that it is a rhombus with diagonals of $6$ and $8$. The area is $\dfrac{1}{2}\times 6\times8$, or $\boxed{\mathrm{(D)}\ 24}$




\textbf{Solution 2}

You can also assign $x$ and $y$ to be $0$. Then you can easily see that the diagonals are $6$ and $8$. Multiply and divide by $2$ to get D. $24$




\textbf{Solution 3}

The graph is symmetric with respect to both coordinate axes, and in the first quadrant it coincides with the graph of the line $3x + 4y = 12.$ Therefore the region is a rhombus, and the area is

 
 \[\text{Area} = 4\left(\frac{1}{2}(4\cdot 3)\right) = 24 \rightarrow \boxed{D}\]

                    \begin{figure}[H]
                        \centering
                        \includegraphics[width=0.5\textwidth]{Solution/AMC_12_12B_2005_7_Solution_Figure_2.png}
                    \end{figure}
                    



	\item (\textit{AMC 12 12B 2005}) For how many values of $a$ is it true that the line $y = x + a$ passes through the
vertex of the parabola $y = x^2 + a^2$ ?


 $\mathrm{(A)}\ 0      \qquad \mathrm{(B)}\ 1      \qquad \mathrm{(C)}\ 2      \qquad \mathrm{(D)}\ 10      \qquad \mathrm{(E)}\ \text{infinitely many}$

 \textbf{Answer: }C



\textbf{Solution}

We see that the vertex of the quadratic function $y = x^2 + a^2$ is $(0,\,a^2)$. The y-intercept of the line $y = x + a$ is $(0,\,a)$. We want to find the values (if any) such that $a=a^2$. Solving for $a$, the only values that satisfy this are $0$ and $1$, so the answer is $\boxed{\mathrm{(C)}\ 2}$



	\item (\textit{AMC 12 12B 2005}) On a certain math exam, $10\%$ of the students got $70$ points, $25\%$ got $80$ points, $20\%$ got $85$ points, $15\%$ got $90$ points, and the rest got $95$ points.  What is the difference between the mean and the median score on this exam?


 $\mathrm{(A)}\ {{{0}}} \qquad \mathrm{(B)}\ {{{1}}} \qquad \mathrm{(C)}\ {{{2}}} \qquad \mathrm{(D)}\ {{{4}}} \qquad \mathrm{(E)}\ {{{5}}}$

 \textbf{Answer: }B



\textbf{Solution}

To begin, we see that the remaining $30\%$ of the students got $95$ points. Assume that there are $20$ students; we see that $2$ students got $70$ points, $5$ students got $80$ points, $4$ students got $85$ points, $3$ students got $90$ points, and $6$ students got $95$ points. The median is $85$, since the $10^{\text{th}}$ and $11^{\text{th}}$ terms are both $85$. The mean is $\dfrac{70\,(2)+80\,(5)+85\,(4)+90\,(3)+95\,(6)}{20}=\dfrac{1720}{20}=86$. The difference between the mean and median, therefore, is $\boxed{\textbf{(B)}\ 1}$.

 



\textbf{Solution 2}

The remaining $30\%$ of the students got $95$ points.The mean is equal to $10\%\cdot70 + 25\%\cdot80 + 20\%\cdot85 + 15\%\cdot90 + 30\%\cdot95 = 86$.The score greater than or equal to $50\%$ of other scores is the median. Since $35\%$ scored $80$ or lower and the next $20\%$ scored $85$, the median is $85$. The difference between the mean and the median is $\boxed{\textbf{(B)}\ 1}$.

 


	\item (\textit{AMC 12 12B 2005}) The first term of a sequence is $2005$. Each succeeding term is the sum of the cubes of the digits of the previous terms. What is the $2005^{\text{th}}$ term of the sequence?


 $\mathrm{(A)}\ {{{29}}} \qquad \mathrm{(B)}\ {{{55}}} \qquad \mathrm{(C)}\ {{{85}}} \qquad \mathrm{(D)}\ {{{133}}} \qquad \mathrm{(E)}\ {{{250}}}$

 \textbf{Answer: }E



\textbf{Solution}

Performing this operation several times yields the results of $133$ for the second term, $55$ for the third term, and $250$ for the fourth term. The sum of the cubes of the digits of $250$ equal $133$, a complete cycle. The cycle is, excluding the first term, the $2^{\text{nd}}$, $3^{\text{rd}}$, and $4^{\text{th}}$ terms will equal $133$, $55$, and $250$, following the fourth term. Any term number that is equivalent to $1\ (\text{mod}\ 3)$ will produce a result of $250$. It just so happens that $2005\equiv 1\ (\text{mod}\ 3)$, which leads us to the answer of $\boxed{\textbf{(E) } 250}$.

 


	\item (\textit{AMC 12 12B 2005}) An envelope contains eight bills: $2$ ones, $2$ fives, $2$ tens, and $2$ twenties.  Two bills are drawn at random without replacement.  What is the probability that their sum is $$20$ or more?


 $\mathrm{(A)}\ {{{\frac{1}{4}}}} \qquad \mathrm{(B)}\ {{{\frac{2}{7}}}} \qquad \mathrm{(C)}\ {{{\frac{3}{7}}}} \qquad \mathrm{(D)}\ {{{\frac{1}{2}}}} \qquad \mathrm{(E)}\ {{{\frac{2}{3}}}}$

 \textbf{Answer: }D



\textbf{Solution 1}

The only way to get a total of \$$20$ or more is if you pick a twenty and another bill, or if you pick both tens. There are a total of $\dbinom{8}{2}=\dfrac{8\times7}{2\times1}=28$ ways to choose $2$ bills out of $8$. There are $12$ ways to choose a twenty and some other non-twenty bill. There is $1$ way to choose both twenties, and also $1$ way to choose both tens. Adding these up, we find that there are a total of $14$ ways to attain a sum of $20$ or greater, so there is a total probability of $\dfrac{14}{28}=\boxed{\textbf{(D) }\frac{1}{2}}$.

 



\textbf{Solution 2}

Another way to do this problem is to use \href{https://artofproblemsolving.com/wiki/index.php?title=Complementary\_counting}{complementary counting}, i.e. how many ways that the sum is less than $20$. Now, you do not have to consider the $2$ twenties, so you have $6$ bills left. $\dbinom{6}{2} = \dfrac{6\times5}{2\times1} = 15$ ways. However, you counted the case when you have $2$ tens, so you need to subtract 1, and you get $14$. Finding the ways to get $20$ or higher, you subtract $14$ from $28$ and get $14$. So the answer is $\dfrac{14}{28} = \boxed{\textbf{(D) }\dfrac{1}{2}}$




\textbf{Solution 3}

There are two cases that work, namely getting at least $1$ twenty, or getting $2$ tens.

 Case $1$: $P(\text{Get at least one twenty}) = 1-P(\text{Do not get a single twenty})=1- \frac{\binom{6}{2}}{\binom{8}{2}}=\frac{28-15}{28}=\frac{13}{28}$

 Case $2$ : $P(\text{Get two tens}) = \frac{1}{\binom{8}{2}} = \frac{1}{28}$

 Summing up our cases, we have $\frac{13}{28}+\frac{1}{28}=\frac{14}{28}=\boxed{\textbf{(D) } \dfrac{1}{2}}$




\textbf{Solution 4}

Note that if a twenty is drawn, anything else that is drawn will create a total greater than $20$; The probability of a twenty being drawn first is $\frac{1}{4}.$ The same could be said for drawing anything, and then drawing a twenty. However, we can only draw something that isn't a twenty first (since we've already accounted for the probability of drawing two twenties). 

 The probability of drawing a non-twenty first, then a twenty second is $\frac{3}{4}\cdot\frac{2}{7}=\frac{3}{14}.$ Finally, we can draw two tens. The probability of this occuring is $\frac{1}{4}\cdot\frac{1}{7}=\frac{1}{28}.$ 

 Adding these three probabilities gives us $\frac{1}{4}+\frac{3}{14}+\frac{1}{28}=\boxed{\textbf{(D) } \dfrac{1}{2}}$

 -Benedict T (countmath1) (edited by AMC\_8)

 



\textbf{Solution 5 (Quick if no time)}

We see that there are $2$ of each bill. We can simplify the question to only drawing $1$ bill out of $4$, and trying to draw $\$ 10$. (Originally you need to draw $\$ 20$, but the bill is halved, so you need to draw only $\$ 10$.) We see that we have $1$ of each bill $\$1$, $\$5$, $\$10$, and $\$20$.Thus, we can easily see the solution is $\boxed{(D) \frac{1}{2}}$.

 


	\item (\textit{AMC 12 12B 2005}) The \href{https://artofproblemsolving.com/wiki/index.php?title=Quadratic\_equation}{quadratic equation} $x^2+mx+n$ has roots twice those of $x^2+px+m$, and none of $m,n,$ and $p$ is zero. What is the value of $n/p$?


 $\mathrm{(A)}\ {{{1}}} \qquad \mathrm{(B)}\ {{{2}}} \qquad \mathrm{(C)}\ {{{4}}} \qquad \mathrm{(D)}\ {{{8}}} \qquad \mathrm{(E)}\ {{{16}}}$

 \textbf{Answer: }D



\textbf{Solution 1}

Let $x^2 + px + m = 0$ have roots $a$ and $b$. Then 

 
 \[x^2 + px + m = (x-a)(x-b) = x^2 - (a+b)x + ab,\]

 so $p = -(a+b)$ and $m = ab$. Also, $x^2 + mx + n = 0$ has roots $2a$ and $2b$, so

 
 \[x^2 + mx + n = (x-2a)(x-2b) = x^2 - 2(a+b)x + 4ab,\]

 and $m = -2(a+b)$ and $n = 4ab$. Thus $\frac{n}{p} = \frac{4ab}{-(a+b)} = \frac{4m}{\frac{m}{2}} = \boxed{\textbf{(D) }8}$. 

 To test that this actually works, consider the quadratics $x^2 + 8x + 16 = 0,\ x^2 + 16x + 64 = 0$.

 



\textbf{Solution 2}

If the roots of $x^2 + mx + n = 0$ are $2a$ and $2b$ and the roots of $x^2 + px + m = 0$ are $a$ and $b$, then using Vieta's formulas,
 \[2a + 2b = -m\]
 \[a + b = -p\]
 \[2a(2b) = n\]
 \[a(b) = m\]Therefore, substituting the second equation into the first equation gives 
 \[m = 2(p)\]and substituting the fourth equation into the third equation gives
 \[n = 4(m)\]Therefore, $n = 8p$, so $\frac{n}{p}= \boxed{\textbf{(D) }8}$



	\item (\textit{AMC 12 12B 2005}) Suppose that $4^{x_1}=5$, $5^{x_2}=6$, $6^{x_3}=7$, ... , $127^{x_{124}}=128$.  What is $x_1x_2...x_{124}$?


 $\mathrm{(A)}\ {{{2}}} \qquad \mathrm{(B)}\ {{{\frac{5}{2}}}} \qquad \mathrm{(C)}\ {{{3}}} \qquad \mathrm{(D)}\ {{{\frac{7}{2}}}} \qquad \mathrm{(E)}\ {{{4}}}$

 \textbf{Answer: }D



\textbf{Solution}

We see that we can re-write $4^{x_1}=5$, $5^{x_2}=6$, $6^{x_3}=7$, ... , $127^{x_{124}}=128$ as $\left(...\left(\left(\left(4^{x_1}\right)^{x_2}\right)^{x_3}\right)...\right)^{x_{124}}=128$ by using substitution. By using the properties of exponents, we know that $4^{x_1x_2...x_{124}}=128$. 

 $4^{x_1x_2...x_{124}}=128\\2^{2x_1x_2...x_{124}}=2^7\\2x_1x_2...x_{124}=7\\x_1x_2...x_{124}=\dfrac{7}{2}$

 Therefore, the answer is $\boxed{\mathrm{(D)}\,\dfrac{7}{2}}$

 \textbf{Alternate Solution} Changing $4^{x_1}=5$ to logarithmic form, we get ${x_1}=\log_{4}{5}$.  We can rewrite this as ${x_1}=\dfrac{\log{5}}{\log{4}}$.  Applying this to the rest, we get $x_1x_2...x_{124}=\dfrac{\log5}{\log4}\cdot\dfrac{\log6}{\log5}\cdot...\cdot\dfrac{\log128}{\log127}=\dfrac{\log5\cdot\log6\cdot...\cdot\log128}{\log4\cdot\log5\cdot...\cdot\log127}=\dfrac{\log128}{\log4}=\log_{4}{128}=\log_{4}{2^7}=7\cdot\log_{4}{2}=7\cdot\dfrac{1}{2}=\boxed{\mathrm{(D)}\,\dfrac{7}{2}}$



	\item (\textit{AMC 12 12B 2005}) A circle having center $(0,k)$, with $k>6$, is tangent to the lines $y=x$, $y=-x$ and $y=6$. What is the radius of this circle?


 $\mathrm{(A)}\ 6\sqrt{2}-6 \qquad \mathrm{(B)}\ 6 \qquad \mathrm{(C)}\ 6\sqrt{2} \qquad \mathrm{(D)}\ 12 \qquad \mathrm{(E)}\ 6+6\sqrt{2}$

 \textbf{Answer: }E



\textbf{Solution}

Let $R$ be the radius of the circle. Draw the two radii that meet the points of tangency to the lines $y = \pm x$. We can see that a square is formed by the origin, two tangency points, and the center of the circle. The side lengths of this square are $R$ and the diagonal is $k = R+6$. The diagonal of a square is $\sqrt{2}$ times the side length. Therefore, $R+6 = R\sqrt{2} \Rightarrow R = \dfrac{6}{\sqrt{2}-1} = 6+6\sqrt{2} \Rightarrow \boxed{\mathrm{E}}$. 

 
                    \begin{figure}[H]
                        \centering
                        \includegraphics[width=0.5\textwidth]{Solution/AMC_12_12B_2005_14_Solution_Figure_1.png}
                    \end{figure}
                    



	\item (\textit{AMC 12 12B 2005}) The sum of four two-digit numbers is $221$. None of the eight digits is $0$ and no two of them are the same. Which of the following is \textbf{not} included among the eight digits?


 $\mathrm{(A)}\ 1      \qquad \mathrm{(B)}\ 2      \qquad \mathrm{(C)}\ 3      \qquad \mathrm{(D)}\ 4      \qquad \mathrm{(E)}\ 5$

 \textbf{Answer: }D



\textbf{Solution 1}

$221$ can be written as the sum of four two-digit numbers, let's say $\overline{ae}$, $\overline{bf}$, $\overline{cg}$, and $\overline{dh}$. Then $221= 10(a+b+c+d)+(e+f+g+h)$. The last digit of $221$ is $1$, and $10(a+b+c+d)$ won't affect the units digits, so $(e+f+g+h)$ must end with $1$. The smallest value $(e+f+g+h)$ can have is $(1+2+3+4)=10$, and the greatest value is $(6+7+8+9)=30$. Therefore, $(e+f+g+h)$ must equal $11$ or $21$. 

 Case 1: $(e+f+g+h)=11$

 The only distinct positive integers that can add up to $11$ is $(1+2+3+5)$. So, $a$,$b$,$c$, and $d$ must include four of the five numbers $(4,6,7,8,9)$. We have $10(a+b+c+d)=221-11=210$, or $a+b+c+d=21$. We can add all of $4+6+7+8+9=34$, and try subtracting one number to get to $21$, but none of them work. Therefore, $(e+f+g+h)$ cannot add up to $11$.

 Case 2: $(e+f+g+h)=21$

 Checking all the values for $e$,$f$,$g$,and $h$ each individually may be time-consuming, instead of only having $1$ solution like Case 1. We can try a different approach by looking at $(a+b+c+d)$ first. If $(e+f+g+h)=21$, $10(a+b+c+d)=221-21=200$, or $(a+b+c+d)=20$. That means $(a+b+c+d)+(e+f+g+h)=21+20=41$. We know $(1+2+3+4+5+6+7+8+9)=45$, so the missing digit is $45-41=\boxed{\mathrm{(D)}\ 4}$




\textbf{Solution 2}

Alternatively, we know that a number is congruent to the sum of its digits mod 9, so $221 \equiv 5 \equiv 1 + 2 + 3 + 4 + 5 + 6 + 7 + 8 + 9 - d \equiv -d$, where $d$ is some digit. Clearly, $\boxed{d = 4}$.

 



\textbf{Solution 3 (Solution 1 without cases)}

Like solution $1$, note that the variables can be expressed as:

 
 \[10(a + b + c + d) + (e + f + g + h) = 221.\]

 The largest possible sum for 4 digits is $(6 + 7 + 8 + 9) = 30$, so the tens digit is under $30$. Therefore, the only possible splits are

 
 \[10(20) + 21 = 221 \quad \text{and} \quad 10(21) + 11 = 221.\]

 Since the sum of the digits $1\text{–}9$ is $45$, removing the largest digit $(9)$  yields a minimum sum of $36$.

 As $21 + 11 = 32 < 36$, the $10(21) + 11$ split is impossible. Therefore, the missing digit is $45 - (20 + 21) = \boxed{d = 4}$.

 


	\item (\textit{AMC 12 12B 2005}) Eight spheres of radius 1, one per octant, are each tangent to the coordinate planes. What is the radius of the smallest sphere, centered at the origin, that contains these eight spheres?


 $\mathrm {(A)}\ \sqrt{2}  \qquad \mathrm {(B)}\ \sqrt{3}  \qquad \mathrm {(C)}\ 1+\sqrt{2}\qquad \mathrm {(D)}\ 1+\sqrt{3}\qquad \mathrm {(E)}\ 3$

 \textbf{Answer: }D



\textbf{Solution}

The eight spheres are formed by shifting spheres of radius $1$ and center $(0, 0, 0)$ $\pm 1$ in the $x, y, z$ directions. Hence, the centers of the spheres are $(\pm 1, \pm 1, \pm 1)$. For a sphere centered at the origin to contain all eight spheres, its radius must be greater than or equal to the longest distance from the origin to one of these spheres. This length is the sum of the distance from $(\pm 1, \pm 1, \pm 1)$ to the origin and the radius of the spheres, or $\sqrt{3} + 1$. To verify this is the longest length, we can see from the triangle inequality that the length from the origin to any other point on the spheres is strictly smaller. Thus, the answer is $\boxed{\mathrm{D}}$.

 



\textbf{Solution 2}

Simplify the problem - in $2$ dimensions, place $4$ circles in each quadrant in the coordinate plane. To find the circle enclosing these $4$ circles, notice that if you connect the $4$ centers as a square, the diameter of the large circle will be the diagonal of the square plus two small radii. 

 Similarly, in $3$ dimensions, consider the cube created by connecting the $8$ centers of the spheres. The diameter of the large spheres will just be the space diagonal of this cube, plus the two radiuses on either side. The side length of the cube is $2$, so the space diagonal will be $2\sqrt{3}$. Then the diameter of the large sphere is $2\sqrt{3} + 2$ and the radius is $\sqrt{3} + 1$ or $\boxed{\mathrm{D}}$.

 


	\item (\textit{AMC 12 12B 2005}) How many distinct four-tuples $(a, b, c, d)$ of rational numbers are there with


 $a \cdot \log_{10} 2+b \cdot \log_{10} 3 +c \cdot \log_{10} 5 + d \cdot \log_{10} 7 = 2005$?


 $\mathrm{(A)}\ 0      \qquad \mathrm{(B)}\ 1      \qquad \mathrm{(C)}\ 17     \qquad \mathrm{(D)}\ 2004   \qquad \mathrm{(E)}\ \text{infinitely many}$

 \textbf{Answer: }B



\textbf{Solution}

Using the laws of \href{https://artofproblemsolving.com/wiki/index.php?title=Logarithms}{logarithms}, the given equation becomes

 
 \[\log_{10}2^{a}+\log_{10}3^{b}+\log_{10}5^{c}+\log_{10}7^{d}=2005\]
 \[\Rightarrow \log_{10}{2^{a}\cdot 3^{b}\cdot 5^{c}\cdot 7^{d}}=2005\]
 \[\Rightarrow 2^{a}\cdot 3^{b}\cdot 5^{c}\cdot 7^{d} = 10^{2005}\]

 As $a,b,c,d$ must all be rational, and there are no powers of $3$ or $7$ in $10^{2005}$, $b=d=0$. Then $2^{a}\cdot 5^{c}=2^{2005}\cdot 5^{2005} \Rightarrow a=c=2005$.

 Only the four-tuple $(2005,0,2005,0)$ satisfies the equation, so the answer is $\boxed{1} \Rightarrow \mathrm{(B)}$.

 


	\item (\textit{AMC 12 12B 2005}) Let $A(2,2)$ and $B(7,7)$ be points in the plane. Define $R$ as the region in the first quadrant consisting of those points $C$ such that $\triangle ABC$ is an acute triangle. What is the closest integer to the area of the region $R$?


 $\mathrm{(A)}\ 25     \qquad \mathrm{(B)}\ 39     \qquad \mathrm{(C)}\ 51     \qquad \mathrm{(D)}\ 60     \qquad \mathrm{(E)}\ 80     \qquad$

 \textbf{Answer: }C



\textbf{Solution}


                    \begin{figure}[H]
                        \centering
                        \includegraphics[width=0.5\textwidth]{Solution/AMC_12_12B_2005_18_Solution_Figure_1.png}
                    \end{figure}
                    

 For angle $A$ and $B$ to be acute, $C$ must be between the two lines that are perpendicular to $\overline{AB}$ and contain points $A$ and $B$. For angle $C$ to be acute, first draw a $45-45-90$ triangle with $\overline{AB}$ as the hypotenuse. Note $C$ cannot be inside this triangle's circumscribed circle or else $\angle C > 90^\circ$. Hence, the area of $R$ is the area of the large triangle minus the area of the small triangle minus the area of the circle, which is $\frac{14^2}{2}-\frac{4^2}{2}-(\frac{5\sqrt{2}}{2})^2\pi=98-8-\frac{25\pi}{2}$, which is approximately $51$. The answer is $\boxed{\mathrm{C}}$.

 Note: When doing this problem, it is important to double check that the circle with diameter $AB$ lies fully in the first quadrant, which luckily it does.

 


	\item (\textit{AMC 12 12B 2005}) Let $x$ and $y$ be two-digit integers such that $y$ is obtained by reversing the digits of $x$. The integers $x$ and $y$ satisfy $x^{2}-y^{2}=m^{2}$ for some positive integer $m$. What is $x+y+m$?


 $\mathrm{(A)}\ 88    \qquad \mathrm{(B)}\ 112   \qquad \mathrm{(C)}\ 116   \qquad \mathrm{(D)}\ 144   \qquad \mathrm{(E)}\ 154   \qquad$

 \textbf{Answer: }E



\textbf{Solution}

let $x=10a+b$, then $y=10b+a$ where $a$ and $b$ are nonzero digits.

 By \href{https://artofproblemsolving.com/wiki/index.php?title=Difference\_of\_squares}{difference of squares},
 \[x^2-y^2=(x+y)(x-y)\]
 \[=(10a+b+10b+a)(10a+b-10b-a)\]
 \[=(11(a+b))(9(a-b))\]

 For this product to be a square, the factor of $11$ must be repeated in either $(a+b)$ or $(a-b)$, and given the constraints it has to be $(a+b)=11$. The factor of $9$ is already a square and can be ignored. Now $(a-b)$ must be another square, and since $a$ cannot be $10$ or greater then $(a-b)$ must equal $4$ or $1$. If $a-b=4$ then $(a+b)+(a-b)=11+4$, $2a=15$, $a=15/2$, which is not a digit. Hence the only possible value for $a-b$ is $1$. Now we have $(a+b)+(a-b)=11+1$, $2a=12$, $a=6$, then $b=5$, $x=65$, $y=56$, $m=33$, and $x+y+m=154\Rightarrow\boxed{\mathrm{E}}$



	\item (\textit{AMC 12 12B 2005}) Let $a,b,c,d,e,f,g$ and $h$ be distinct elements in the set


 
 \[\{-7,-5,-3,-2,2,4,6,13\}.\]


 What is the minimum possible value of


 
 \[(a+b+c+d)^{2}+(e+f+g+h)^{2}?\]
$\mathrm{(A)}\ 30     \qquad \mathrm{(B)}\ 32     \qquad \mathrm{(C)}\ 34     \qquad \mathrm{(D)}\ 40     \qquad \mathrm{(E)}\ 50$

 \textbf{Answer: }C



\textbf{Solution}

The sum of the set is $-7-5-3-2+2+4+6+13=8$, so if we could have the sum in each set of parenthesis be $4$ then the minimum value would be $2(4^2)=32$. Considering the set of four terms containing $13$, this sum could only be even if it had zero, two, or four odd terms. If it had all four odd terms then it would be $13-7-5-3=-2$, and with two odd terms then its minimum value is $13-7+2-2=6$, so we cannot achieve two sums of $4$. The closest we could have to $4$ and $4$ is $3$ and $5$, which can be achieved through $13-7-5+2$ and $6-3-2+4$. So the minimum possible value is $3^2+5^2=34\Rightarrow\boxed{\mathrm{C}}$.

 



\textbf{Solution 2 (better explanation of 1st)}

Trying out the values for a bit leads to us getting $34$ with $(-3, -2, 2, 6)$ in 1 set for a total of 9 and the other numbers giving a total of 25. 

 We start off with trying to get a $2(4^2) = 32$ solution, so we only need to find one set with a sum of 4 (the other will automatically have a sum of 4). 

 We can add 7 to each number. We are now trying to get a set with 32.

 
 \[(0, 2, 4, 5, 9, 11, 13, 20)\]

 Observe that there are 4 odd and 4 even values. To get 32 from 4 numbers $a,b,c,d$, the numbers either have to be 4 odds, 4 evens, or 2 odds and evens. The 4 odd numbers do not work (38) nor the 4 even numbers (26). 

 If $a,b,c,d$ has 2 odd numbers and 2 even numbers we can divide them into two cases: whether $a,b,c,d$ has 20 or doesn't have 20.

 Case 1: $a,b,c,d$ doesn't have 20.

 The maximum sum of the 2 even numbers are 2 + 4 = 6 and the maximum sum of the 2 odd numbers are 11 + 13 = 24. This is not enough to reach the 32 we need.

 Case 2: $a,b,c,d$ has 20.

 The minimum sum of the 2 even numbers are 20 + 0 = 20 and the minimum sum of the 2 odd numbers are 5 + 9 = 14. This is already greater than the 32 we need.

 We have proven that we cannot partition the 8 numbers into two groups of 4 with the same sum, so the smallest value is $\boxed{C}$.

 



\textbf{Solution 3}

Apply Cauchy's inequality to bound the sum, yielding $(1+1)\left((a+b+c+d)^{2}+(e+f+g+h)^{2}\right) \ge (a+b+c+d+e+f+g+h)^2$ and rearrange to get $(a+b+c+d)^{2}+(e+f+g+h)^{2} \ge 32$, where equality is attained iff $\frac{a+b+c+d}{1}=\frac{e+f+g+h}{1}=4$, which is impossible from the numbers given upon observing. Therefore, we can eliminate choices $\mathrm{(A)}$ and $\mathrm{(B)}$. Now experiment to get that $13-7-5+2$ and $6-3-2+4$ give $34$, which must be the minimum since we've eliminated the smaller answer choices. $\implies \boxed{\mathrm{C}}$



	\item (\textit{AMC 12 12B 2005}) A positive integer $n$ has $60$ divisors and $7n$ has $80$ divisors.  What is the greatest integer $k$ such that $7^k$ divides $n$?


 $\mathrm{(A)}\ {{{0}}} \qquad \mathrm{(B)}\ {{{1}}} \qquad \mathrm{(C)}\ {{{2}}} \qquad \mathrm{(D)}\ {{{3}}} \qquad \mathrm{(E)}\ {{{4}}}$

 \textbf{Answer: }C



\textbf{Solution}

We may let $n = 7^k \cdot m$, where $m$ is not divisible by 7. Using the fact that the number of divisors function $d(n)$ is multiplicative, we have $d(n) = d(7^k)d(m) = (k+1)d(m) = 60$. Also, $d(7n) = d(7^{k+1})d(m) = (k+2)d(m) = 80$. These numbers are in the ratio 3:4, so $\frac{k+1}{k+2} = \frac{3}{4} \implies k = 2 \Rightarrow \boxed{\mathrm{C}}$.

 


	\item (\textit{AMC 12 12B 2005}) A sequence of complex numbers $z_{0}, z_{1}, z_{2}, ...$ is defined by the rule


 
 \[z_{n+1} = \frac {iz_{n}}{\overline {z_{n}}},\]


 where $\overline {z_{n}}$ is the \href{https://artofproblemsolving.com/wiki/index.php?title=Complex\_conjugate}{complex conjugate} of $z_{n}$ and $i^{2}=-1$. Suppose that $|z_{0}|=1$ and $z_{2005}=1$. How many possible values are there for $z_{0}$?


 $\textbf{(A)}\ 1 \qquad  \textbf{(B)}\ 2 \qquad  \textbf{(C)}\ 4 \qquad  \textbf{(D)}\ 2005 \qquad  \textbf{(E)}\ 2^{2005}$

 \textbf{Answer: }E



\textbf{Solution 1}

Since $|z_0|=1$, let $z_0=e^{i\theta_0}$, where $\theta_0$ is an \href{https://artofproblemsolving.com/wiki/index.php?title=Argument}{argument} of $z_0$. We will prove by induction that $z_n=e^{i\theta_n}$, where $\theta_n=2^n(\theta_0+\frac{\pi}{2})-\frac{\pi}{2}$.

 Base Case: trivial

 Inductive Step: Suppose the formula is correct for $z_k$, then
 \[z_{k+1}=\frac{iz_k}{\overline {z_k}}=i e^{i\theta_k} e^{i\theta_k}=e^{i(2\theta_k+\pi/2)}\]Since
 \[2\theta_k+\frac{\pi}{2}=2\cdot 2^n(\theta_0+\frac{\pi}{2})-\pi+\frac{\pi}{2}=2^{n+1}(\theta_0+\frac{\pi}{2})-\frac{\pi}{2}=\theta_{n+1}\]the formula is proven

 $z_{2005}=1\Rightarrow \theta_{2005}=2k\pi$, where $k$ is an integer. Therefore,
 \[2^{2005}(\theta_0+\frac{\pi}{2})=(2k+\frac{1}{2})\pi\]
 \[\theta_0=\frac{k}{2^{2004}}\pi+\left(\frac{1}{2^{2006}}-\frac{1}{2}\right)\pi\]The value of $\theta_0$ only matters \href{https://artofproblemsolving.com/wiki/index.php?title=Modulo}{modulo} $2\pi$. Since $\frac{k+2^{2005}}{2^{2004}}\pi\equiv\frac{k}{2^{2004}}\pi\mod 2\pi$, k can take values from 0 to $2^{2005}-1$, so the answer is $2^{2005}\Rightarrow\boxed{\mathrm{E}}$




\textbf{Solution 2}

Let $z_0 = \cos \theta + i\sin \theta = e^{i\theta}$.
 \[z_1 = \frac {iz_{0}}{\overline {z_{0}}} = \frac{i(\cos \theta + i\sin \theta)}{\cos \theta - i\sin \theta} = i(\cos \theta + i\sin \theta)^2 = i(\cos 2\theta + i\sin 2\theta) = ie^{i2\theta}\]
 \[z_2 = \frac {iz_{1}}{\overline {z_{1}}} = \frac {i(ie^{i2\theta})}{-ie^{-i2\theta}} = -ie^{i4\theta}\]
 \[z_3 = \frac {iz_{2}}{\overline {z_{2}}} = \frac {i(-ie^{i4\theta})}{ie^{-i4\theta}} = -ie^{i8\theta}\]

 Repeating through this recursive process, we can quickly see that for $n>1$
 \[z_n = -i \cdot e^{i2^n\theta}\]Plugging in $n=2005$, we get 
 \[z_{2005} = -ie^{i2^{2005}\theta} = 1\]Rearranging terms, we see that $e^{i2^{2005}\theta}=i$.This means that $e^{i2^{2005}\theta}$ lies at point $i$ on the unit circle in the complex plane, which happens when $2^{2005}\theta = \frac{\pi}{2} + 2k\pi$. Solving for $\theta$, we get $\theta=\frac{\pi(2k+\frac{1}{2})}{2^{2005}}$ where $k = 0,1,2...2^{2005}-1$. So the answer is $2^{2005}\Rightarrow\boxed{\mathrm{E}}$. (Author: Patrick Yin)

 



\textbf{Solution 3}

Note that for any complex number $z$, we have $|z|=|\overline z|$. Therefore, the magnitude of $\frac{iz_n}{|z_n|}$ is always $1$, meaning that all of the numbers in the sequence $z_k$ are of magnitude $1$.

 Another property of complex numbers is that $z\overline z=|z|^2$. For the numbers in our sequence, this means $z\overline z=1$, so $\overline z=z^{-1}$. Rewriting our recursive condition with these facts, we now have
 \[z_{n+1}=\frac{iz_n}{z_n^{-1}}=iz_n^2.\]Solving for $z_n$ here, we obtain
 \[z_n=\frac{\pm\sqrt{z_{n+1}}}i=-i\cdot(\pm\sqrt{z_{n+1}}).\]It is seen that there are two values of $z_n$ which correspond to one value of $z_{n+1}$. That means that there are two possible values of $z_{2004}$, four possible values of $z_{2003}$, and so on. Therefore, there are $2^{2005}$ possible values of $z_0$, giving the answer as $\boxed{\mathrm{(E)}\text{ }2^{2005}}$.

 



\textbf{Solution 4 (more accurate solution 2)}

Let $z_0=e^{i\theta}$. Then $z_1=\frac{iz_0}{\overline{z_0}}=\frac{ie^{i\theta}}{e^{-i\theta}}=ie^{i2\theta}$, $z_2=\frac{iz_1}{\overline{z_1}}-e^{i2^2 \theta}$, $z_3=\frac{iz_2}{\overline{z_2}}=-ie^{i2^3 \theta}$, $z_4=\frac{iz_3}{\overline{z_3}}=e^{i2^4\theta}$. Now we see that every for every positive integer $4k$, $z_{4k}=e^{i2^{4k}\theta}$ so $z_{2004}=e^{i2^{2004}\theta}$ and $z_{2005}=ie^{i2^{2005}\theta}=1 \iff e^{i(2^{2005}\theta + \frac{\pi}{2})}=1$, which has $2^{2005}$ solutions of the form 
 \[\theta=\frac{\frac{3\pi}{2}+2\pi k}{2^{2005}}\] for $k \in \{0, 1, \dots 2^{2005}-1\}$. $\implies \boxed{\mathrm{E}}$



	\item (\textit{AMC 12 12B 2005}) Let $S$ be the set of ordered triples $(x,y,z)$ of real numbers for which


 
 \[\log_{10}(x+y) = z \text{ and } \log_{10}(x^{2}+y^{2}) = z+1.\]
There are real numbers $a$ and $b$ such that for all ordered triples $(x,y,z)$ in $S$ we have $x^{3}+y^{3}=a \cdot 10^{3z} + b \cdot 10^{2z}.$ What is the value of $a+b?$
$\textbf{(A)}\ \frac {15}{2} \qquad  \textbf{(B)}\ \frac {29}{2} \qquad  \textbf{(C)}\ 15 \qquad  \textbf{(D)}\ \frac {39}{2} \qquad  \textbf{(E)}\ 24$

 \textbf{Answer: }B



\textbf{Solution 1}

Let $x + y = s$ and $x^2 + y^2 = t$.   Then, $\log(s)=z$ implies $\log(10s) = z+1= \log(t)$,so $t=10s$.   Therefore, $x^3 + y^3 = s\cdot\dfrac{3t-s^2}{2} = s(15s-\dfrac{s^2}{2})$.  Since $s = 10^z$, we find that $x^3 + y^3 = 15\cdot10^{2z} - (1/2)\cdot10^{3z}$.  Thus, $a+b = \frac{29}{2}$ $\Rightarrow$ $\boxed{B}$




\textbf{Solution 2}

First, remember that $x^3 + y^3$ factors to 

 


	\item (\textit{AMC 12 12B 2005}) All three vertices of an equilateral triangle are on the parabola $y=x^2$, and one of its sides has a slope of $2$.  The $x$-coordinates of the three vertices have a sum of $m/n$, where $m$ and $n$ are relatively prime positive integers.  What is the value of $m+n$?


 $\mathrm{(A)}\ {{{14}}} \qquad \mathrm{(B)}\ {{{15}}} \qquad \mathrm{(C)}\ {{{16}}} \qquad \mathrm{(D)}\ {{{17}}} \qquad \mathrm{(E)}\ {{{18}}}$

 \textbf{Answer: }A



\textbf{Solution 1 (Complex numbers)}

Let the three points be at $A = (x_1, x_1^2)$, $B = (x_2, x_2^2)$, and $C = (x_3, x_3^2)$, such that the slope between the first two is $2$, and $A$ is the point with the least $y$-coordinate.

 Then, we have $\textrm{Slope of }AC = \frac{x_1^2 - x_3^2}{x_1 - x_3} = x_1 + x_3$. Similarly, the slope of $BC$ is $x_2 + x_3$, and the slope of $AB$ is $x_1 + x_2 = 2$. The desired sum is $x_1 + x_2 + x_3$, which is equal to the sum of the slopes divided by $2$.

 To find the slope of $AC$, we note that it comes at a $60^{\circ}$ angle with $AB$. Thus, we can find the slope of $AC$ by multiplying the two complex numbers $1 + 2i$ and $1 + \sqrt{3}i$. What this does is generate the complex number that is at a $60^{\circ}$ angle with the complex number $1 + 2i$. Then, we can find the slope of the line between this new complex number and the origin:
 \[(1+2i)(1+\sqrt{3}i)\]
 \[= 1 - 2\sqrt{3} + 2i + \sqrt{3}i\]
 \[\textrm{Slope } = \frac{2 + \sqrt{3}}{1 - 2\sqrt{3}}\]
 \[= \frac{2 + \sqrt{3}}{1 - 2\sqrt{3}} \cdot \frac{1 + 2\sqrt{3}}{1 + 2\sqrt{3}}\]
 \[= \frac{8 + 5\sqrt{3}}{-11}\]
 \[= \frac{-8 - 5\sqrt{3}}{11}.\]The slope $BC$ can also be solved similarly, noting that it makes a $120^{\circ}$ angle with $AB$:
 \[(1+2i)(-1+\sqrt{3}i)\]
 \[= -1 - 2\sqrt{3} - 2i + \sqrt{3}i\]
 \[\textrm{Slope } = \frac{\sqrt{3} - 2}{-2\sqrt{3} - 1}\]
 \[= \frac{2 - \sqrt{3}}{1 + 2\sqrt{3}} \cdot \frac{1 - 2\sqrt{3}}{1 - 2\sqrt{3}}\]

 At this point, we start to notice a pattern: This expression is equal to $\frac{2 + \sqrt{3}}{1 - 2\sqrt{3}} \cdot \frac{1 + 2\sqrt{3}}{1 + 2\sqrt{3}}$, except the numerators of the first fractions are conjugates! Notice that this means that when we multiply out, the rational term will stay the same, but the coefficient of $\sqrt{3}$ will have its sign switched. This means that the two complex numbers are conjugates, so their irrational terms cancel out.

 Our sum is simply $2 - 2\cdot\frac{8}{11} = \frac{6}{11}$, and thus we can divide by $2$ to obtain $\frac{3}{11}$, which gives the answer $\boxed{\mathrm{(A)}\ 14}$.

 



\textbf{Solution 2}


                    \begin{figure}[H]
                        \centering
                        \includegraphics[width=0.5\textwidth]{Solution/AMC_12_12B_2005_24_Solution_Figure_1.png}
                    \end{figure}
                    

 Using the slope formula and differences of squares, we find: 

 $a+b$ = the slope of $AB$, 

 $b+c$ = the slope of $BC$, 

 $a+c$ = the slope of $AC$. 

 So the value that we need to find is the sum of the slopes of the three sides of the triangle divided by $2$. Without loss of generality, let $AB$ be the side that has the smallest angle with the positive $x$-axis. Let $J$ be an arbitrary point with the coordinates $(1, 0)$. Translate the triangle so $A$ is at the origin. Then $\tan(BOJ) = 2$. Since the slope of a line is equal to the tangent of the angle formed by the line and the positive x- axis, the answer is $\dfrac{\tan(BOJ) + \tan(BOJ+60) + \tan(BOJ-60)}{2}$. 

 Using $\tan(BOJ) = 2$, and the tangent addition formula, this simplifies to $\dfrac{3}{11}$, so the answer is $3 + 11 = \boxed{\mathrm{(A)}\ 14}$



	\item (\textit{AMC 12 12B 2005}) Six ants simultaneously stand on the six \href{https://artofproblemsolving.com/wiki/index.php?title=Vertex}{vertices} of a regular \href{https://artofproblemsolving.com/wiki/index.php?title=Octahedron}{octahedron}, with each ant at a different vertex. Simultaneously and independently, each ant moves from its vertex to one of the four adjacent vertices, each with equal \href{https://artofproblemsolving.com/wiki/index.php?title=Probability}{probability}. What is the probability that no two ants arrive at the same vertex?


 $\mathrm{(A)}\ \frac {5}{256} \qquad\mathrm{(B)}\ \frac {21}{1024} \qquad\mathrm{(C)}\ \frac {11}{512} \qquad\mathrm{(D)}\ \frac {23}{1024} \qquad\mathrm{(E)}\ \frac {3}{128}$

 \textbf{Answer: }A



\textbf{Solution 1}

We approach this problem by counting the number of ways ants can do their desired migration, and then multiply this number by the probability that each case occurs.

 Let the octahedron be $ABCDEF$, with points $B,C,D,E$ \href{https://artofproblemsolving.com/wiki/index.php?title=Coplanar}{coplanar}. Then the ant from $A$ and the ant from $F$ must move to plane $BCDE$. Suppose, without loss of generality, that the ant from $A$ moved to point $B$. Then, we must consider three cases.

 
            \begin{itemize}
        \item Case 1: Ant from point $F$ moved to point $C$
            \end{itemize}
            On the plane, points $B$ and $C$ are taken. The ant that moves to $D$ can come from either $E$ or $C$. The ant that moves to $E$ can come from either $B$ or $D$. Once these two ants are fixed, the other two ants must migrate to the "poles" of the octahedron, points $A$ and $F$. Thus, there are two degrees of freedom in deciding which ant moves to $D$, two degrees of freedom in deciding which ant moves to $E$, and two degrees of freedom in deciding which ant moves to $A$. Hence, there are $2 \times 2 \times 2=8$ ways the ants can move to different points.

 
            \begin{itemize}
        \item Case 2: Ant from point $F$ moved to point $D$
            \end{itemize}
            On the plane, points $B$ and $D$ are taken. The ant that moves to $C$ must be from $B$ or $D$, but the ant that moves to $E$ must also be from $B$ or $D$. The other two ants, originating from points $C$ and $E$, must move to the poles. Therefore, there are two degrees of freedom in deciding which ant moves to $C$ and two degrees of freedom in choosing which ant moves to $A$. Hence, there are $2 \times 2=4$ ways the ants can move to different points.

 
            \begin{itemize}
        \item Case 3: Ant from point $F$ moved to point $E$
            \end{itemize}
            By symmetry to Case 1, there are $8$ ways the ants can move to different points.

 Given a point $B$, there is a total of $8+4+8=20$ ways the ants can move to different points. We oriented the square so that point $B$ was defined as the point to which the ant from point $A$ moved. Since the ant from point $A$ can actually move to four different points, there is a total of $4 \times 20=80$ ways the ants can move to different points.

 Each ant acts independently, having four different points to choose from. Hence, each ant has probability $1/4$ of moving to the desired location. Since there are six ants, the probability of each case occuring is $\frac{1}{4^6} = \frac{1}{4096}$. Thus, the desired answer is $\frac{80}{4096}= \boxed{\frac{5}{256}} \Rightarrow \mathrm{(A)}$.

 



\textbf{Solution 2}

Let $f(n)$ be the number of \href{https://artofproblemsolving.com/wiki/index.php?title=Cycle}{cycles} of length $n$ the can be walked among the vertices of an octahedron. For example, $f(3)$ would represent the number of ways in which an ant could navigate $2$ vertices and then return back to the original spot. Since an ant cannot stay still, $f(1) = 0$. We also easily see that $f(2) = 1, f(3) = 2$. 

 Now consider any four vertices of the octahedron. All four vertices will be connected by edges except for one pair. Let’s think of this as a \href{https://artofproblemsolving.com/wiki/index.php?title=Square}{square} with one diagonal (from top left to bottom right). EDIT: This part is wrong as if you choose the 4 vertices that have a cross section as a square, there exists no connecting diagonal.  

 
                    \begin{figure}[H]
                        \centering
                        \includegraphics[width=0.5\textwidth]{Solution/AMC_12_12B_2005_25_Solution_Figure_1.png}
                    \end{figure}
                    

 Suppose an ant moved across this diagonal; then the ant at the other end can only move across the diagonal (which creates 2-cycle, bad) or it can move to another vertex, but then the ant at that vertex must move to the spot of the original ant (which creates 3-cycle, bad). Thus none of the ants can navigate the diagonal and can either shift clockwise or counterclockwise, and so $f(4) = 2$.

 For $f(6)$, consider an ant at the top of the octahedron. It has four choices. Afterwards, it can either travel directly to the bottom, and then it has $2$ ways back up, or it can travel along the sides and then go to the bottom, of which simple counting gives us $6$ ways back up. Hence, this totals $4 \times (2+6) = 32$.

 Now, the number of possible ways is given by the sum of all possible cycles,
 \[a \cdot f(2) \cdot f(2) \cdot f(2) + b \cdot f(2) \cdot f(4) + c \cdot f(3) \cdot f(3) + d \cdot f(6)\]

 where the \href{https://artofproblemsolving.com/wiki/index.php?title=Coefficient}{coefficients} represent the number of ways we can configure these cycles. To find $a$, fix any face, there are $4$ adjacent faces to select from to complete the cycle. From the four remaining faces there are only $2$ ways to create cycles, hence $a = 8$.

 To find $b$, each cycle of $2$ faces is distinguished by their common edge, and there are $12$ edges, so $b = 12$.

 To find $c$, each three-cycle is distinguished by the vertex, and there are $8$ edges. However, since the two three-cycles are indistinguishable, $c = 8/2 = 4$.

 Clearly $d = 1$. Finally,

 
 \[8(1)(1)(1) + 12(1)(2) + 4(2)(2) + (32) = 80\]

 Each bug has $4$ possibilities to choose from, so the probability is $\frac{80}{4^6} = \frac{5}{256}$.

 



\textbf{Solution 3}

We use the idea of cycles. 

 Case 1: Two 3-cycles.We have 4 ways to divide the faces into cycles. There are 8 faces, and once a face has been chosen, the other three vertices must form the other cycle of length 3. Then there are $2^2$ ways to choose the direction of the cycle.  There are $4 \cdot 4 = 16.$ 

 Case 2: Three 2-cycles.There are $4 \cdot 2 = 8$ ways to do so. 

 Case 3: One 4-cycle, one 2-cycle.There are $12 \cdot 2 = 24$ ways here because there are 12 different edges for the 2-cycle and we can choose the direction of the 4-cycle. 

 Case 4: One 6-cycle.Visualizing the octahedron as a 4-sided pyramid pointing up above one pointing down, we look at the ways an ant can start from the top vertex and visit all vertices without revisiting one along the way and returning to the top. Because the first step must be from the top to the middle square ring, we choose one vertex to move to and will multiply the final result by four.

 We divide the paths into two sub-cases.

 Sub-case 1: The ant continues to the bottom vertex.  In this case, the ant has visited three corners of a square, but can not next visit the fourth corner or there will be no way to connect the remaining two vertices of the octahedron.  Thus, the ant must next visit one of the other two vertices, and that selection decides the remainder of the path.  Thus, this sub-case has 2 possible paths.

 Sub-case 2: The ant continues along the "equator" square ring. By symmetry, there are two choices.  Choose one, and we will multiply the final result by two.  From that second point on the equator, there are two choices.  If it continues to a third point on the equator, there is only one path to complete the cycle.  If it next moves to the bottom vertex, there are two ways to complete the cycle. This gives a total of three possible paths.  Multiplying by 2 gives 6 for this sub-case.

 There are $(2 + 6) \cdot 4 = 32$ ways here. 

 Overall, there are $80$ cases.  Answer is $80/2^{12} = 5/256$




\textbf{Solution 4}

Let $ACEDFB$ be the octahedron such that points $CDEF$ are coplanar, $C/D$ are opposite each other and $E/F$ are opposite each other.

 The set of points that the ants from $A/B$ can travel to is $\{C,D,E,F\}$, the set of points that the ants from $C/D$ can travel to is $\{A,B,E,F\}$ and the set of points that the ants from $E/F$ can travel to is $\{A,B,C,D\}$. So the problem is analogous to finding the number of ways to take 2 elements from each of these sets to make the set $\{A,B,C,D,E,F\}$.

 
            \begin{itemize}
        \item Case 1: Ants from $E/F$ travel to $A/B$
            \end{itemize}
            If this occurs then the ants from $C/D$ must travel to $E/F$ and therefore the ants from $E/F$ must travel to $C/D$. So this gives 1 case.

 
            \begin{itemize}
        \item Case 2: Ants from $E/F$ travel to $A/C$
            \end{itemize}
            If this happens then one of the ants from $C/D$ must go to $B$ and one of the ants from $A/B$ must go to $D$ and the remaining ant from $C/D$ can choose to go to $E$ or $F$ and then the remaining ant from $A/B$ has only 1 choice. So this gives 2 cases.

 The ants from $E/F$ traveling to $A/B$ is equivalent to them traveling to $C/D$ (both remove 2 elements from another set). And the ants from $E/F$ traveling to $A/C$ is equivalent to them traveling to $A/D$ or $B/C$ or $B/D$ (both remove 1 element from the other 2 sets). Also, for each of the 3 pairs of ants there are 2 ways the pair of ants can go to 2 points (for example, $E\mapsto A$ and $F\mapsto B$ or $E\mapsto B$ and $F\mapsto A$) so there are $2^{3}(2\cdot1 + 4\cdot2) = 80$ possibilties. There are a total of $2^{12}$ ways the ants can move so the answer is $\frac{80}{2^{12}} = \boxed{\textbf{(A) } \frac{5}{256}}$.~LuisFonseca123

 


\end{enumerate}